\clearpage
\thispagestyle{empty}

\begin{center}
	\vspace*{4cm}
	{\Huge \textbf{KẾT THÚC PHẦN 2}} \\[1.5cm]
	
	\rule{0.8\textwidth}{0.4pt} \\[0.8cm]
	
	\parbox{0.9\textwidth}{
		\centering
		\textit{
			Phần 2 đã đưa bạn vào một hành trình thực chiến, biến những lý thuyết phức tạp 
			thành các kỹ năng và sản phẩm cụ thể. Từ những bước đầu tiên trong việc 
			xử lý dữ liệu đến việc huấn luyện và triển khai các mô hình ngôn ngữ lớn, 
			bạn đã được trang bị một bộ công cụ toàn diện để giải quyết các bài toán 
			NLP trong thế giới thực. Phần này không chỉ dạy bạn cách sử dụng các thư viện, 
			mà còn rèn luyện một tư duy MLOps bền vững. Với kiến thức từ cả hai phần, 
			bạn đã có đủ hành trang để tự tin bước đi trên con đường của một chuyên gia NLP.
		}
	} \\[1cm]
	
	\rule{0.8\textwidth}{0.4pt} \\[2cm]
\end{center}

\noindent
\textbf{Tóm lược nội dung Phần 2:}
\begin{itemize}[leftmargin=*]
	\item \textbf{Chương 1 – Quy trình làm việc và Công cụ Xử lý Dữ liệu:} 
	Phác thảo vòng đời dự án NLP tinh gọn, các kỹ thuật thu thập dữ liệu (APIs, Web Scraping), 
	làm sạch và gán nhãn (Pandas, Polars, Label Studio, Snorkel, LLM-based), 
	tăng cường dữ liệu, và quản lý phiên bản với DVC.
	
	\item \textbf{Chương 2 – Huấn luyện và Đánh giá Mô hình:}
	Làm chủ hệ sinh thái Hugging Face (`transformers`, `datasets`, `accelerate`, `evaluate`),
	xây dựng các hàm đánh giá thực tế (`compute\_metrics`), theo dõi thí nghiệm
	(Weights \& Biases, MLflow), huấn luyện phân tán với DeepSpeed,
	toàn bộ chuỗi hậu huấn luyện với TRL (SFT/QLoRA, DPO, GRPO),
	và đánh giá tái lập được với LM Evaluation Harness.
	
	\item \textbf{Chương 3 – Tối ưu hóa, Triển khai và Vận hành (MLOps):}
	Tìm hiểu về Cơ sở dữ liệu Vector, các thư viện tối ưu hóa (`bitsandbytes`, `peft`, AWQ/GPTQ, ONNX),
	đóng gói và phục vụ mô hình (FastAPI, Docker), các framework chuyên dụng (vLLM với prefix caching
	và giải mã suy đoán, SGLang, BentoML, Triton), và vòng đời MLOps hoàn chỉnh (giám sát, tái huấn luyện).
	
	\item \textbf{Chương 4 – Công thức Xây dựng các Ứng dụng Phổ biến (Recipes):} 
	Hướng dẫn từng bước để xây dựng các ứng dụng thực tế: Tìm kiếm Ngữ nghĩa, 
	Hệ thống RAG, Fine-tuning cho Phân loại và NER, xây dựng Tác tử đơn giản, 
	và Trích xuất Thông tin có Cấu trúc.
\end{itemize}